\documentclass[12pt,letterpaper]{article}
\usepackage[utf8]{inputenc}
\usepackage[spanish]{babel}
\usepackage[margin=2.5cm]{geometry}
\usepackage{xcolor}
\usepackage{hyperref}

\begin{document}

\begin{center}
    \Large \textbf{Curso de Seguridad Informática} \\
    \large \textbf{Suite de Auditoría en Python} \\
    \vspace{0.3cm}
    \normalsize Prof. César Rodríguez \\
\end{center}

\vspace{0.5cm}

\noindent \textbf{Fecha:} \today \\
\textbf{Asunto:} Tareas y plazos para la Fase II (Enumeración y Ataques) \\
\textbf{Plazo de entrega:} Semanas 5 a 7

\vspace{0.5cm}

\noindent Estimados estudiantes,

\vspace{0.3cm}

Espero que se encuentren muy bien. Tras haber concluido exitosamente la etapa inicial de nuestro proyecto, me dirijo a ustedes para anunciar el inicio oficial de la \textbf{Fase II: Enumeración y Ataques}. 

Durante las \textbf{semanas 5, 6 y 7} de nuestro calendario, cada grupo se enfocará en desarrollar un módulo ofensivo o de enumeración profunda basado en los capítulos de nuestro libro de texto:

\begin{itemize}
    \item \textbf{Grupo 1:} Enumeración y Banner Grabbing (\texttt{banner\_grabber.py})
    \item \textbf{Grupo 2:} Enumeración NetBIOS/SMB (\texttt{smb\_enumerator.py})
    \item \textbf{Grupo 3:} Fuerza Bruta FTP (\texttt{bruteforce\_ftp.py})
    \item \textbf{Grupo 4:} Fuerza Bruta Web e Integración General (\texttt{bruteforce\_web.py})
\end{itemize}

\vspace{0.3cm}
\noindent \textbf{Instrucciones y Entregables:}
\begin{enumerate}
    \item \textbf{Material Teórico:} Se ha publicado en el repositorio el documento \textit{Guía de Desarrollo: Parte II} (\texttt{Guia\_Fase\_II.pdf}). Es obligatorio que cada grupo lea la teoría que le corresponde y comprenda el flujo de trabajo antes de empezar a programar.
    \item \textbf{Desarrollo:} Las plantillas de código orientadas a objetos ya se encuentran disponibles en la carpeta \texttt{modulos/} del repositorio. 
    \item \textbf{Entrega (Pull Request):} El código finalizado deberá ser enviado mediante un \textbf{Pull Request} en GitHub antes de finalizar la \textbf{Semana 7}. Si necesitan recordar el proceso, consulten la \textit{Guía para la Creación de Pull Requests} en la carpeta \texttt{docs/}.
\end{enumerate}

\vspace{0.3cm}
\noindent \textbf{Recordatorios Críticos de Arquitectura:}
\begin{itemize}
    \item \textbf{Carpetas:} Todos los desarrollos deben guardarse y ejecutarse estrictamente dentro de la carpeta \texttt{modulos/}. No creen directorios adicionales.
    \item \textbf{Regla de Oro:} Se mantiene la prohibición absoluta de modificar o alterar los archivos asignados a otros grupos.
    \item \textbf{Independencia:} Gracias a la Programación Orientada a Objetos, ningún grupo necesita esperar a otro para avanzar en su código.
\end{itemize}

Quedo a su entera disposición en caso de dudas teóricas o bloqueos técnicos. ¡Mucho éxito en sus desarrollos, confío plenamente en las habilidades que han demostrado hasta ahora!

\vspace{1cm}
\noindent Atentamente, \\
\vspace{0.5cm} \\
\textbf{Prof. César Rodríguez}

\end{document}